% =============================================================
\section{Method}\label{sec:method}
% =============================================================

We adopt the notation of \cref{tab:notation}. The pipeline is:
\begingroup
\setlength{\abovedisplayskip}{4pt}
\setlength{\belowdisplayskip}{4pt}
\begin{equation}
  I_e(t) \xrightarrow{\text{field}} \boldsymbol{\Phi}
  \xrightarrow{G\,\cdot} \mathbf{f}
  \xrightarrow{\text{encode}} \mathbf{I}_{\text{ecs}}
  \xrightarrow{\text{integrate}} \mathbf{v}(t).
  \label{eq:pipeline}
\end{equation}
\endgroup

\begin{table}[t]
\centering
\small
\caption{Notation. Lengths in $\mu$m, currents in $\mu$A or nA, voltages
in mV, time in ms, conductivity in S/m.}
\label{tab:notation}
\begin{tabular}{@{}lll@{}}
\toprule
Symbol & Meaning & Units \\ \midrule
$N$ & Number of compartments & --- \\
$T$ & Number of timesteps & --- \\
$\mathbf{x}_j,\,\mathbf{x}_e$ & Compartment center, electrode position & $\mu$m \\
$d_{je}$ & $\lVert \mathbf{x}_j - \mathbf{x}_e \rVert$ & $\mu$m \\
$I_e(t)$ & Electrode current (cathodic negative) & $\mu$A \\
$\sigma$ & Extracellular conductivity & S/m \\
$\phi_j(t),\,\boldsymbol{\Phi}$ & Extracellular potential, $\boldsymbol{\Phi}_{jt}=\phi_j(t)$ & mV \\
$v_j(t),\,\mathbf{v}$ & Transmembrane voltage & mV \\
$G$ & Axial voltage diffusion operator & $1/$ms \\
$c_j,\,A_j$ & Specific capacitance, surface area & $\mu$F/cm$^2$,\,$\mu$m$^2$ \\
$f_j(t) = [G\boldsymbol{\phi}(t)]_j$ & Activating function at compartment $j$ & mV/ms \\
\bottomrule
\end{tabular}
\end{table}

\paragraph{Extracellular potential from point sources.}
For an electrode $e$ delivering $I_e(t)$ in a homogeneous, isotropic,
infinite volume conductor of conductivity $\sigma$, superposition of
monopoles gives
\[
  \phi_j(t) \;=\; \sum_e \frac{I_e(t)}{4\pi\sigma\,d_{je}}.
\]
Defining the transfer factor $h_{je} = 10^3 / (4\pi\sigma\,d_{je})$
(in $\mathrm{mV}/\mu\mathrm{A}$, with the $10^3$ absorbing the unit
conversion from $\mu$A and $\mu$m to mV) factorises space and time, so
that for a single electrode
$\boldsymbol{\Phi} = \mathbf{h}\,\mathbf{I}^\top \in \mathbb{R}^{N\times T}$.

\paragraph{The discrete activating function.}
Jaxley solves the spatially discretised cable equation
\[
  c_j A_j \frac{dv_j}{dt}
  = \sum_{k\sim j} g_{jk}(v_k - v_j) - A_j I_{\text{ion}}(v_j,\mathbf{s}_j) + i_{\text{ext},j},
\]
where $g_{jk}$ is the raw axial conductance (units S) between
compartments $j$ and $k$, $c_j$ is specific membrane capacitance, and
$A_j$ is compartment surface area. Internally, Jaxley folds $c_j A_j$
and unit conversions into a single voltage diffusion operator $G$ with
entries $\tilde g_{jk} = g_{jk} / (c_j A_j)$ (in $1/$ms; note the
asymmetric row-$j$ normalisation, which is why $G$ is symmetric only
when $c_j A_j$ is uniform across compartments), so that
$\dot{\mathbf{v}} = G\mathbf{v} + (\text{membrane terms})$. The sign
pattern of $G$ is the weighted graph Laplacian:
$G_{jk} = +\tilde g_{jk}$ for adjacent $k\neq j$,
$G_{jj} = -\sum_{k\neq j} \tilde g_{jk}$, and zero otherwise. Rows sum
to zero (current conservation).

An extracellular field $\phi_j(t)$ shifts the axial driving force from
$v_k - v_j$ to $(v_k+\phi_k)-(v_j+\phi_j)$; by linearity of $G$, this
is equivalent to adding $G\boldsymbol{\phi}$ as a forcing term, the
discrete analogue of Rattay's continuous activating function
$\rho_a^{-1}\partial^2\phi/\partial x^2$~\cite{rattay1986analysis}.
On a uniform straight cable with constant radius and compartment
length $\Delta x$, $[G\boldsymbol{\phi}]_j$ at interior compartments
reduces to a constant multiple of the standard
$[1,-2,1]/\Delta x^2$ stencil for $\partial^2\phi/\partial x^2$, with
$O(\Delta x^2)$ truncation error (verified in
\cref{sec:results-verify}).
On branched morphologies, we delegate the elimination of the
branchpoint pseudo-nodes that Jaxley introduces at cable junctions to
its own solver routine, keeping the operator in sync with upstream
solver changes (\cref{app:branchpoint-elimination}).

\paragraph{Equivalent injected current.}
The activating function $G\boldsymbol{\phi}$ has units mV/ms (a voltage
rate), but Jaxley's stimulus API accepts current in nA. We encode the
forcing as a per-compartment current $\mathbf{I}_{\text{ecs}}$ chosen
so that, after Jaxley's internal current-to-voltage-rate conversion,
the integrator sees exactly $G\boldsymbol{\phi}$. The unit-conversion
algebra is given in \cref{app:equivalent-current}; the membrane
capacitance and area factors that appear in the encoding cancel
exactly in the round trip and are an artifact of working through the
public API rather than modifying the solver.

\paragraph{Implementation.}
The pipeline of \cref{eq:pipeline} is realised as four small JAX-native
modules wired together by an experiment object.
\texttt{field.point\_source\_potential} computes $\boldsymbol{\Phi}$ from
electrode positions, currents, and $\sigma$, with a $1\,\mu$m floor on
$d_{je}$ to regularise the $1/d$ singularity when an electrode lands inside
a compartment.
\texttt{discretization.build\_voltage\_operator\_G} returns $G$ by
delegating to Jaxley's \texttt{\_compute\_transition\_matrix} (the same
routine the simulator uses for its backward-Euler step) and stripping
branchpoint pseudo-nodes.
\texttt{equivalent\_current.phi\_e\_to\_ecs\_nA} applies the encoding
above.
An \texttt{ECSExperiment} class precomputes $G$, $\mathbf{c}$, $\mathbf{A}$,
and the compartment centres in its constructor, and exposes
\texttt{simulate\_waveform} for a single integration and
\texttt{simulate\_and\_extract} for one integration plus response-feature
extraction.

All operations are JAX-traceable: \texttt{jit}, \texttt{vmap},
\texttt{grad}, and \texttt{jax.sharding.NamedSharding} compose with
the pipeline without modification, so a one-dimensional device mesh is
enough to data-parallelise the batch axis across a TPU pod-slice. Two
implementation choices matter at scale and unlock the full-fidelity
throughput regime of \cref{sec:res-throughput}: $G$ is stored as a
sparse \texttt{BCOO} matrix to avoid the $(B,\,N,\,N)$ XLA broadcast
that dense $G$ triggers under \texttt{vmap}, and
\texttt{simulate\_waveform} forwards a \texttt{checkpoint\_lengths}
factorisation of $T$ into \texttt{jx.integrate} so the backward-Euler
scan rematerialises in $\sim\!\sqrt{T}$ chunks rather than holding the
full $(T,\,B,\,N)$ voltage history in HBM.

